\subsection{Definition of the essential upper integral}

DEFINITION 1. --- For every function \(f \in \mathcal{F}_{+}(\mathrm{T})\) one calls essential upper integral of \(f\) with respect to \(\mu\) , and denotes by \(\mu^{\bullet}(f)\) , the supremum, finite or not, of the set of numbers \(\mu^{*}(f\varphi_{\mathrm{K}})\) , where \(\mathrm{K}\) runs over the set of compact subsets of \(\mathrm{T}\) . For every subset \(\mathrm{A}\) of \(\mathrm{T}\) , one sets \(\mu^{\bullet}(\mathrm{A}) = \mu^{\bullet}(\varphi_{\mathrm{A}})\) .

The notations \(\int^{\bullet}f d\mu ,\int^{\bullet}f(t)d\mu (t),\int^{\bullet}f\mu\) are also used.

Since \(f\varphi_{\mathrm{K}} \leqslant f\) for every compact subset \(\mathrm{K}\) of \(\mathrm{T}\) , one has

\[
\int^ {\bullet} f d \mu \leqslant \int^ {*} f d \mu .\tag{1}
\]

It can happen that \(\mu^{\bullet}(f) \neq \mu^{*}(f)\) ; for, the condition \(\mu^{*}(f) = 0\) means that \(f\) is negligible, whereas the condition \(\mu^{\bullet}(f) = 0\) means that \(f\) is locally negligible (Ch. IV, §5, No.~2, Prop. 5), and there may exist locally negligible sets that are not negligible (Ch. IV, §1, Exer. 5).

The mapping \(\mu^{\bullet}\) of \(\mathcal{F}_{+}(\mathrm{T})\) into \(\overline{\mathbf{R}}\) coincides with \(\mu\) on \(\mathcal{H}_{+}(\mathrm{T})\) . It follows that two measures \(\mu_{1}\) and \(\mu_{2}\) such that \(\mu_{1}^{\bullet} = \mu_{2}^{\bullet}\) are equal.

PROPOSITION 1. --- a) If \(f\) and \(g\) are two numerical functions \(\geqslant 0\) that are equal locally almost everywhere, then \(\mu^{\bullet}(f) = \mu^{\bullet}(g)\) .

\begin{enumerate}
\def\labelenumi{\alph{enumi})}
\setcounter{enumi}{1}
\item
  If \(f\) and \(g\) are two numerical functions \(\geqslant 0\) such that \(f \leqslant g\) , then \(\mu^{\bullet}(f) \leqslant \mu^{\bullet}(g)\) .
\item
  If \(f\) is a numerical function \(\geqslant 0\) , and \(\alpha\) is a number \(\geqslant 0\) , then \(\mu^{\bullet}(\alpha f) = \alpha \mu^{\bullet}(f)\) .
\item
  If \(f\) and \(g\) are two numerical functions \(\geqslant 0\) , then \(\mu^{\bullet}(f + g) \leqslant \mu^{\bullet}(f) + \mu^{\bullet}(g)\) .
\item
  If \((f_n)_{n\in \mathbf{N}}\) is an increasing sequence of numerical functions \(\geqslant 0\) , and if \(f = \lim_{n\to \infty}f_n\) , then \(\mu^{\bullet}(f) = \lim_{n\to \infty}\mu^{\bullet}(f_n)\) .
\end{enumerate}

The properties a), b), c), d) may be deduced immediately from the corresponding properties of the upper integral: a) from Proposition 6 of Ch. IV, §2, No.~3 and Proposition 5 of Ch. IV, §5, No.~2; b), c), d) from Propositions 10, 11, 12 of Ch. IV, §1, No.~3. To establish e), denote by K the set of compact subsets of T; by Theorem 3 of Ch. IV, §1, No.~3, we have

\[
\begin{array}{c} \lim _ {n \to \infty} \mu^ {\bullet} (f _ {n}) = \sup _ {n \in \mathbf {N}} \sup _ {K \in \mathfrak {K}} \mu^ {*} (f _ {n} \varphi_ {K}) = \sup _ {K \in \mathfrak {K}} \sup _ {n \in \mathbf {N}} \mu^ {*} (f _ {n} \varphi_ {K}) \\ = \sup _ {K \in \mathfrak {K}} \mu^ {*} (f \varphi_ {K}) = \mu^ {\bullet} (f)  . \end{array}
\]

Equality holds in the relation d) if f and g are measurable, by Cor. 4 of Th. 5, Ch. IV, §5, No.~6. More generally, we have the following result:

PROPOSITION 2. --- Let \(f,g,h\) be three elements of \(\mathcal{F}_{+}\) ; if \(g\) and \(h\) are measurable, then

\[
\int^ {\bullet} f (g + h) d \mu = \int^ {\bullet} f g d \mu + \int^ {\bullet} f h d \mu .\tag{2}
\]

One is immediately reduced to the proof of the analogous formula for the upper integral. Since \(f(g + h) = fg + fh\) (with the convention that \(0 \cdot (+\infty) = 0\) ), we have

\[
\int^ {*} f (g + h) d \mu \leqslant \int^ {*} f g d \mu + \int^ {*} f h d \mu ;
\]

it remains to establish the reverse inequality. Let \(u\) be a lower semi-continuous function such that \(u \geqslant f(g + h)\) . Set \(v = \frac{u}{g + h}\) on the set where \(g + h > 0\) , and \(v = +\infty\) on the set where \(g + h = 0\) ; then \(v \geqslant f\) and \(u \geqslant v(g + h)\) , whence

\[
\int^ {*} v (g + h) d \mu \leqslant \int^ {*} u d \mu
\]

and consequently, v being measurable (Ch. IV, §5, No.~6, Cor. 4 of Th. 5),

\[
\begin{array}{r l} \int^ {*} f g d \mu + \int^ {*} f h d \mu & \leqslant \int^ {*} v g d \mu + \int^ {*} v h d \mu \\ & = \int^ {*} v (g + h) d \mu \leqslant \int^ {*} u d \mu , \end{array}
\]

which implies the desired inequality since u is arbitrary.

COROLLARY. --- Let \(f\) be a function \(\geqslant 0\) , \((g_n)\) a sequence of measurable functions \(\geqslant 0\) ; then \(\int^{\bullet} f\left(\sum_{n} g_n\right) d\mu = \sum_{n} \left(\int^{\bullet} fg_n d\mu\right)\) .

For the case of a finite sequence, this is an immediate consequence of Prop. 2. The case of an infinite sequence may be deduced from this by means of Prop. 1, e).

PROPOSITION 3. --- For every finite number \(\alpha \geqslant 0\) and every pair of measures \(\mu, \nu\) on T,

\[
(\alpha \mu) ^ {\bullet} = \alpha \mu^ {\bullet}
\]

\[
(\mu + \nu) ^ {\bullet} = \mu^ {\bullet} + \nu^ {\bullet}.
\]

Moreover, the relation \(\mu \leqslant \nu\) implies \(\mu^{\bullet} \leqslant \nu^{\bullet}\) .

The proof is immediate from the analogous statement in Ch. IV (§1, No.~3, Prop. 15).

PROPOSITION 4. --- For every numerical function \(f \geqslant 0\) that is lower semi-continuous on T, \(\mu^{\bullet}(f) = \mu^{*}(f)\) .

For, let g be a function in \(\mathcal{K}_{+}(\mathrm{T})\) such that \(g \leqslant f\) . If K is the (compact) support of g, then \(\mu(g) \leqslant \mu^{*}(f\varphi_{\mathrm{K}}) \leqslant \mu^{\bullet}(f)\) . It follows, by the definition of upper integral, that \(\mu^{*}(f) \leqslant \mu^{\bullet}(f)\) , therefore \(\mu^{*}(f) = \mu^{\bullet}(f)\) (formula (1)).

